\documentclass[12pt]{article}
\usepackage[spanish]{babel}
\usepackage{amsmath}
\usepackage{amssymb}
\usepackage{amsfonts}
\usepackage{array}
\usepackage{geometry}
\usepackage{multirow}
\geometry{a4paper, margin=2.5cm}

\title{Métricas de Evaluación para Clasificación}
\author{Modelado y simulación I}
\date{}

\begin{document}

\maketitle

\section{Matriz de Confusión}

La matriz de confusión es la base para todas las métricas de clasificación. Para un problema de clasificación binaria, se define:

\begin{table}[h]
\centering
\begin{tabular}{c|c|c|c|}
\cline{3-3}
\multicolumn{2}{c|}{} & \multicolumn{2}{c}{Predicción} \\
\cline{3-4}
\multicolumn{2}{c|}{} & Positivo & Negativo \\
\cline{2-4}
\multirow{2}{*}{Real} & Positivo & VP & FN \\
\cline{2-4}
& Negativo & FP & VN \\
\cline{2-4}
\end{tabular}
\caption{Matriz de Confusión para Clasificación Binaria}
\end{table}

donde:
\begin{itemize}
\item $VP$: Verdaderos Positivos (clase positiva correctamente clasificada)
\item $VN$: Verdaderos Negativos (clase negativa correctamente clasificada)
\item $FP$: Falsos Positivos (clase negativa clasificada como positiva)
\item $FN$: Falsos Negativos (clase positiva clasificada como negativa)
\end{itemize}

%\section{Métricas por Clase}

\subsection{Precisión (Precision)}

Mide la proporción de predicciones positivas que son correctas:

\begin{equation}
\text{Precision} = \frac{VP}{VP + FP}
\end{equation}

Interpretación: De todas las muestras que el modelo clasificó como positivas, ¿cuántas realmente lo son?

\subsection{Recall (Sensibilidad)}

Mide la proporción de muestras positivas que fueron correctamente identificadas:

\begin{equation}
\text{Recall} = \frac{VP}{VP + FN}
\end{equation}

Interpretación: De todas las muestras realmente positivas, ¿cuántas detectó el modelo?

\subsection{F1-Score}

Es la media armónica entre precisión y recall:

\begin{equation}
F_1 = 2 \cdot \frac{\text{Precision} \cdot \text{Recall}}{\text{Precision} + \text{Recall}} = \frac{2VP}{2VP + FP + FN}
\end{equation}

La media armónica penaliza valores extremos, dando un balance entre ambas métricas.

\subsection{Soporte}

El soporte es simplemente el número de muestras reales de cada clase:

\begin{equation}
\text{Soporte}_c = \text{Número de muestras de la clase } c
\end{equation}

\section{Promedios Globales}

Para un problema de clasificación con $C$ clases, donde para cada clase $c$ tenemos:
\begin{itemize}
\item $P_c$: Precisión de la clase $c$
\item $R_c$: Recall de la clase $c$
\item $F_c$: F1-Score de la clase $c$
\item $S_c$: Soporte de la clase $c$ (número de muestras reales)
\end{itemize}

\subsection{Promedio Macro (Macro Avg)}

El promedio macro calcula la media aritmética simple de las métricas de cada clase, dando el mismo peso a todas las clases independientemente de su tamaño:

\begin{equation}
\text{Precision}_{\text{macro}} = \frac{1}{C} \sum_{c=1}^{C} P_c
\end{equation}

\begin{equation}
\text{Recall}_{\text{macro}} = \frac{1}{C} \sum_{c=1}^{C} R_c
\end{equation}

\begin{equation}
F_{1,\text{macro}} = \frac{1}{C} \sum_{c=1}^{C} F_c
\end{equation}

\textbf{Propiedades:}
\begin{itemize}
\item Cada clase contribuye con igual peso
\item Útil cuando todas las clases son igualmente importantes
\item No se ve afectado por el desbalance de clases
\end{itemize}

\subsection{Promedio Ponderado (Weighted Avg)}

El promedio ponderado calcula la media de las métricas ponderada por el soporte de cada clase:

\begin{equation}
\text{Precision}_{\text{weighted}} = \frac{\sum_{c=1}^{C} P_c \cdot S_c}{\sum_{c=1}^{C} S_c}
\end{equation}

\begin{equation}
\text{Recall}_{\text{weighted}} = \frac{\sum_{c=1}^{C} R_c \cdot S_c}{\sum_{c=1}^{C} S_c}
\end{equation}

\begin{equation}
F_{1,\text{weighted}} = \frac{\sum_{c=1}^{C} F_c \cdot S_c}{\sum_{c=1}^{C} S_c}
\end{equation}

\textbf{Propiedades:}
\begin{itemize}
\item Las clases con mayor soporte tienen mayor influencia
\item Refleja el rendimiento global del modelo en el dataset
\item Útil cuando se quiere evaluar el desempeño en la distribución real
\end{itemize}

\section{Exactitud (Accuracy)}

La exactitud es la proporción de predicciones correctas sobre el total:

\begin{equation}
\text{Accuracy} = \frac{VP + VN}{VP + VN + FP + FN}
\end{equation}

Para múltiples clases:

\begin{equation}
\text{Accuracy} = \frac{\sum_{c=1}^{C} VP_c}{\sum_{c=1}^{C} (VP_c + FN_c)}
\end{equation}

donde $VP_c$ son los verdaderos positivos de la clase $c$ y $FN_c$ son los falsos negativos.

\textbf{Limitación:} En datasets desbalanceados, la exactitud puede ser engañosa, ya que puede estar dominada por la clase mayoritaria.

\section{Validación Cruzada para Estimación de Precisión}

La precisión del modelo k-NN (y cualquier modelo de machine learning) se estima mediante validación cruzada. Para $k$ folds:

\begin{equation}
\text{Accuracy}_{\text{CV}} = \frac{1}{k} \sum_{i=1}^{k} \text{Accuracy}_i
\end{equation}

donde $\text{Accuracy}_i$ es la exactitud en el fold $i$.

La desviación estándar de las exactitudes proporciona una medida de la variabilidad:

\begin{equation}
\sigma_{\text{CV}} = \sqrt{\frac{1}{k-1} \sum_{i=1}^{k} (\text{Accuracy}_i - \text{Accuracy}_{\text{CV}})^2}
\end{equation}

%\section{Relación entre Métricas}

\subsection{Relación entre F1 y las métricas componentes}

El F1-Score puede expresarse como:

\begin{equation}
F_1 = \frac{2}{\frac{1}{\text{Precision}} + \frac{1}{\text{Recall}}}
\end{equation}

Esta es la media armónica, que es menor o igual que la media aritmética:

\begin{equation}
F_1 \leq \frac{\text{Precision} + \text{Recall}}{2}
\end{equation}

\subsection{Relación entre Accuracy y otras métricas}

La accuracy puede descomponerse como:

\begin{equation}
\text{Accuracy} = \frac{VP + VN}{VP + VN + FP + FN}
\end{equation}

En términos de precision y recall:

\begin{equation}
\text{Accuracy} = \frac{VP + VN}{(VP + VN + FP + FN)}
\end{equation}

No es una función directa de precision y recall, ya que también depende de los verdaderos negativos.

\section{Ejemplo Numérico}

Supongamos una matriz de confusión:

\begin{table}[h]
\centering
\begin{tabular}{c|c|c|c|}
\cline{3-3}
\multicolumn{2}{c|}{} & \multicolumn{2}{c}{Predicción} \\
\cline{3-4}
\multicolumn{2}{c|}{} & Positivo & Negativo \\
\cline{2-4}
\multirow{2}{*}{Real} & Positivo & 80 & 20 \\
\cline{2-4}
& Negativo & 15 & 85 \\
\cline{2-4}
\end{tabular}
\end{table}

\textbf{Métricas para clase positiva:}
\begin{align}
\text{Precision} &= \frac{80}{80 + 15} = \frac{80}{95} \approx 0.8421 \\
\text{Recall} &= \frac{80}{80 + 20} = \frac{80}{100} = 0.8000 \\
F_1 &= 2 \cdot \frac{0.8421 \cdot 0.8000}{0.8421 + 0.8000} \approx 0.8205
\end{align}

\textbf{Métricas para clase negativa:}
\begin{align}
\text{Precision} &= \frac{85}{85 + 20} = \frac{85}{105} \approx 0.8095 \\
\text{Recall} &= \frac{85}{85 + 15} = \frac{85}{100} = 0.8500 \\
F_1 &= 2 \cdot \frac{0.8095 \cdot 0.8500}{0.8095 + 0.8500} \approx 0.8293
\end{align}

\textbf{Promedios:}
\begin{align}
\text{Precision}_{\text{macro}} &= \frac{0.8421 + 0.8095}{2} = 0.8258 \\
\text{Precision}_{\text{weighted}} &= \frac{0.8421 \cdot 100 + 0.8095 \cdot 100}{200} = 0.8258 \\
\text{Accuracy} &= \frac{80 + 85}{200} = \frac{165}{200} = 0.8250
\end{align}



\end{document}